\section{Results}\label{sec:results}

This section is mechanically generated from the latest run of
\texttt{tranche-cli mc --config config/paris\_intermediate.yaml}; every
numerical statement traces back to a row of
\texttt{artifacts/results.csv} or to
\texttt{data/processed/calibrated\_params.yaml}. The Monte Carlo sample
size used in the headline run is $N = \nPaths{}$ paths.

\subsection{Headline pricing}

\Cref{tab:headline} reports the fair price (mean discounted PV across
paths), the resulting ratio to par, the annualised mean realised return,
the realised volatility, the Sharpe ratio, the 95\% Value-at-Risk and the
probability of a negative annualised return for each of the five
instruments under the Gaussian one-factor copula baseline.


\begin{table}[ht]
\centering
\caption{Headline pricing and risk metrics under the Gaussian one-factor copula baseline, without insurance. All monetary values in millions of euros; returns and volatilities annualised.}
\label{tab:headline}
\small
\begin{tabular}{lrrrrrrrr}
\toprule
Instrument & Initial price (\EUR M) & Fair price (\EUR M) & Fair / Par & Mean ann.\ ret.\ & Vol. & Sharpe & VaR$_{95}$ & Pr.\ negative \\
\midrule
Model A & 0.45 & 0.50 & 1.107 & 0.95\% & 1.29\% & -2.16 & 1.18\% & 22.9\% \\
Model B & 0.45 & 0.41 & 0.908 & -1.13\% & 2.01\% & -2.43 & 4.72\% & 70.1\% \\
Equity & 7.88 & 8.11 & 1.030 & -1.15\% & 7.35\% & -0.66 & 10.98\% & 52.2\% \\
Mezzanine & 11.03 & 8.67 & 0.786 & -2.82\% & 3.68\% & -1.78 & 9.50\% & 98.7\% \\
Senior & 12.60 & 11.83 & 0.939 & -0.63\% & 0.21\% & -21.26 & 0.89\% & 100.0\% \\
\bottomrule
\end{tabular}
\end{table}


Three takeaways.
\begin{enumerate}
  \item \textbf{The senior tranche prices at \fairToParSenior{} of par.}
        Its annualised mean return is \meanReturnSenior{} and the
        probability of a negative annualised return is
        \probNegSenior{}. The low realised volatility (a few hundred basis
        points) confirms the structural protection delivered by the
        attachment point at 60\%.
  \item \textbf{Model A (single apartment) and Model B (SAS pool) deliver
        materially different risk distributions} even though their
        \emph{expected} return is broadly similar. Model A shows
        \probNegModela{} of paths with negative realised return ---
        markedly lower than Model B's \probNegModelb{} --- because the
        bimodal distribution (apartment defaults vs survives) keeps a
        majority of paths positive.
  \item \textbf{The equity tranche is the high-variance, high-upside
        instrument.} Its realised return distribution carries a long
        right tail driven by the residual claim on terminal sale
        proceeds; the probability of a negative return is
        \probNegEquity{}.
\end{enumerate}

\Cref{fig:riskreturn-bars,fig:pareto} visualise the same data set as bars
of risk-adjusted ratios and as a risk-return scatter respectively.

\begin{figure}[ht]
  \centering
  \includegraphics[width=\linewidth]{fig_riskreturn_bars.pdf}
  \caption{Risk-adjusted ratios across instruments under the Gaussian
  copula baseline, without insurance. Sharpe, Sortino, Omega and Calmar
  are computed on the annualised return distribution and the mean
  cumulative cash-flow trajectory (see \cref{sec:risk}).}
  \label{fig:riskreturn-bars}
\end{figure}

\begin{figure}[ht]
  \centering
  \includegraphics[width=0.85\linewidth]{fig_pareto_frontier.pdf}
  \caption{Risk-return frontier across the five instruments: realised
  standard deviation of the annualised return on the x-axis, mean
  annualised return on the y-axis. Insurance variants are plotted with
  square markers; the without-insurance baseline uses circles.}
  \label{fig:pareto}
\end{figure}

\subsection{Sensitivity to the choice of credit model}

The same instruments priced under each of the three credit models yield
qualitatively similar but quantitatively different fair-to-par ratios:
\begin{itemize}
  \item Gaussian copula: senior $= \seniorFairGauss{}$.
  \item Student-t copula: senior $= \seniorFairStudent{}$ --- the
        non-zero lower tail dependence (\cref{fig:tail-dep}) shifts more
        loss probability into the senior layer.
  \item Cox doubly stochastic: senior $= \seniorFairCox{}$ ---
        time-varying hazards smooth the cumulative-loss tail relative to
        the static asset-value models, producing a slightly narrower
        distribution (see \cref{fig:loss-dist}).
\end{itemize}

The size of the Gaussian-to-Student gap on the senior tranche
illustrates the Donnelly-Embrechts (2010) and MacKenzie-Spears (2014)
critique: a Gaussian copula is a defensible baseline only when the
joint distribution is thin-tailed. Switching to $\nu = 5$ Student-t on
the Paris calibration drops the senior fair value from
\seniorFairGauss{} to \seniorFairStudent{}.

\subsection{Fair coupons}\label{ssec:fair-coupons}

When the contract coupons are exogenously fixed (3\,\% senior, 6\,\%
mezzanine, 0 equity) they do not necessarily make the PV of each tranche
equal to its par. We therefore solve, separately under each credit model,
the coupon levels that set $\mathrm{PV} = \text{par}$ for the senior and
the mezzanine tranches via Brent's method on the residual function. The
equity tranche stays as a residual claimant. \Cref{tab:fair-coupons}
reports the result.


\begin{table}[ht]
\centering
\caption{Fair coupons that set tranche PV equal to par at $t = 0$, solved by Brent's method for each credit model. Equity stays as the residual claimant.}
\label{tab:fair-coupons}
\small
\begin{tabular}{l r r r r}
\toprule
Tranche & Contract coupon & Cox Intensity & Gaussian Copula & Student T Copula \\
\midrule
Senior & 3.00\% & — & 3.82\% & 3.94\% \\
Equity & 0.00\% & 0.00\% & 0.00\% & 0.00\% \\
\bottomrule
\end{tabular}
\end{table}


Two observations. First, the Brent solver returns a finite senior fair
coupon under the Gaussian and Student-t copulas (3.82\% and 3.94\%
respectively, both above the 3\% contract level) but not under the Cox
model: the time-varying hazard pushes the loss distribution tail
sufficiently for the residual function to have no root in the
$[0, 30\,\%]$ bracket. Second, the Student-t copula's fair senior coupon
exceeds the Gaussian by approximately 12 basis points --- the tail-
dependence premium quantified empirically. The mezzanine tranche has
\emph{no} solution under the baseline cash-flow structure: the available
rent, after maintenance and the senior coupon, is insufficient to bring
mezz to par for any coupon level. We treat this as a finding rather than
a model bug — the structural cash deficiency identifies a candidate
fix (cash trap, overcollateralisation test, structural enhancement)
left as future work.

\subsection{Sensitivity to $\rho$, $\nu$ and $p_d$}\label{ssec:sensitivity}

The fair-to-par metric is most exposed to the copula correlation and the
default-rate calibration. \Cref{fig:price-vs-rho} shows the per-tranche
curve as $\rho$ varies from 0 to 0.45 under the Gaussian copula baseline.
Senior is monotonically decreasing in $\rho$; mezzanine non-monotone;
equity rises slightly as the correlation concentrates loss mass into the
junior layer it absorbs first. The corresponding tornado on the senior
fair-to-par over a $\pm 20\,\%$ shift in $(\rho, p_d, \text{gross yield},
\text{maintenance})$ is shown in \cref{fig:tornado}.

\begin{figure}[ht]
  \centering
  \begin{minipage}{0.48\textwidth}
    \includegraphics[width=\linewidth]{fig_tranche_price_vs_rho.pdf}
    \caption{Per-tranche fair-to-par across a $\rho$ sweep under the
    Gaussian one-factor copula baseline.}
    \label{fig:price-vs-rho}
  \end{minipage}\hfill
  \begin{minipage}{0.48\textwidth}
    \includegraphics[width=\linewidth]{fig_tornado_sensitivity.pdf}
    \caption{Tornado on the senior fair-to-par. Each bar pair shows the
    range over a $\pm 20\,\%$ shock to one parameter; bars are sorted by
    absolute impact.}
    \label{fig:tornado}
  \end{minipage}
\end{figure}

\subsection{Bootstrap confidence intervals}\label{ssec:bootstrap-ci}

Every fair-price estimate carries a 1\,000-resample percentile bootstrap
interval (column \texttt{fair\_price\_ci\_width} in
\texttt{artifacts/results.csv}). \Cref{fig:fanchart} plots the resulting
95\,\% CIs as whiskers on top of the headline bar chart for each of the
three credit models. The senior tranche has the tightest CI (low
realised loss volatility); the equity tranche the widest. The headline
figure (\cref{fig:headline}) collapses everything onto a single chart
suitable as the report's cover.

\begin{figure}[ht]
  \centering
  \includegraphics[width=\linewidth]{fig_fairprice_fanchart.pdf}
  \caption{Bootstrap 95\,\% confidence intervals on fair price per
  instrument and credit model. Whiskers are 1\,000-resample percentile
  bootstrap of the Monte Carlo path-level PVs.}
  \label{fig:fanchart}
\end{figure}

\begin{figure}[ht]
  \centering
  \includegraphics[width=\linewidth]{fig_headline_summary.pdf}
  \caption{Headline summary. Each cluster is one instrument; bars within
  a cluster are the three credit models; whiskers are the 95\,\% CIs
  from \cref{fig:fanchart}.}
  \label{fig:headline}
\end{figure}

\subsection{Insurance}

The actuarial premium charged under the standard-deviation principle is
\insuranceAnnualPremium{} per year. The option-theoretic alternative is
reported in \texttt{results\_meta.json} alongside. Note that the realised
default probability in the Paris-intermediate scenario
(\pdTerminal{} over $T = \horizonYears{}$ years) is calibrated from a
3\% annual rate which is realistic but on the high end; sensitivity to
lower default rates is explored in
\texttt{notebooks/05\_stress\_and\_sensitivity.ipynb}.
